\documentclass[a4paper,11pt]{article}
\usepackage{graphicx}
\usepackage[a4paper,margin=2.5cm]{geometry}
\usepackage{xcolor} 
\usepackage{tikz}    
\usepackage[french,provide=*]{babel}
\usepackage{array} % Inclure le package array
\usepackage[colorlinks=true, linkcolor=black, urlcolor=blue]{hyperref}
\geometry{left=3cm, right=3cm, top=2.5cm, bottom=2.5cm} % Définir les marges

\definecolor{companycolor}{RGB}{0,82,155}  


\date{}

\newcolumntype{C}[1]{>{\centering\arraybackslash}m{#1}}

\begin{document}
\begin{titlepage}
    \centering
    \includegraphics[width=7cm]{../Images/hager.png}\hfill
    \includegraphics[width=7cm]{../Images/uimm.jpg}\par\vspace{1cm} % Remplacez par le logo de votre entreprise
    {\scshape\LARGE HagerGroup}\hspace{5cm}
    {\scshape\LARGE UIMM Alsace\par}
    \vspace{1cm}
    {\scshape\Large Rapport de fin de projet d'études E6 \par}
    \vspace{1.5cm}
    {\huge\bfseries Capteur de Qualité d'Air Intérieur(QAI) \par}
    \vspace{2cm}
    {\Large\itshape Apprenti :}
    \vspace{0.5cm}
    {\Large Timothée Delhelle \par}
    {\Large\itshape Tuteur :}
    \vspace{0.5cm}
    {\Large Nicolas Freyd \par}
    \vfill
    {\Large Diplôme : \par}
    \vspace{0.5cm}
    {\Large BTS CIEL (Cybersécurité, Informatique et Réseaux, Electronique) \par}
    \vfill
    {\Large Promotion : \par}
    \vspace{0.5cm}
    {\Large 2023-2025 \par}
\end{titlepage}

    \tableofcontents
    \newpage

    \section{Introduction}
        Hager Electro est une filiale de Hager Group fabriquant principalement des disjoncteurs au sein de
    plusieurs usines situées majoritairement en Alsace et en Allemagne.
    Actuellement la qualité de l’air des usines n’est pas analysée, ce qui peut représenter un risque
    pour la santé des personnes due à de potentielles particules fines.
    Pour cela, il s’agit d’utiliser un microcontrôleur récupérant les données de capteurs de particules
    fines puis les envoyant à un serveur permettant le traitement, la visualisation et le stockage de ses
    données.
    Actuellement, des serveurs sont en places mais n’ont pas les logiciels requis et aucuns capteurs
    ne sont présents au sein des usines.
    \newpage

    \section{Présentation de l'entreprise}
    \subsection{Localisation}
    \includegraphics[width=420px]{../Images/Localisations.png}
    \subsection{Secteur d'activité}
    \includegraphics[width=420px]{../Images/Marques.png}
    \subsection{Produits phares}
    \includegraphics[width=420px]{../Images/Secteur.png}
    \subsection{Chiffres clés}
    \includegraphics[width=420px]{../Images/Chiffres.png}
    \newpage
    \subsection{Mes Missions}
    \begin{itemize}
        \item Sauvegarde du parc informatique des usines du site d'Obernai à l'aide de la solution Veeam.
        \newline \newline
        \includegraphics[width=420px]{../Images/veeam_backup.png}
        \item Cartographie du réseau à l'aide de l'outil XMind. La cartographie inclut les PC de production, les PLC (automates) ainsi que les périphériques IP tels que des caméras, imprimantes, IHM (Interface Homme Machine) etc.
        \newline 
        \includegraphics[width=420px]{../Images/cartographie.png}
        \newpage
        \item Provisionning de nouveaux PC via PXE à l'aide d'un template disponible via Ivanti. Les nouveaux PC sont provisionnés automatiquement puis enrollés dans l'AD et leurs informations renseignées dans la console Ivanti à l'aide de l'outil Custom Data Form.
        \newline
        \includegraphics[width=420px]{../Images/provisionning.png}
        \item Migration de PC depuis Windows 7 vers Windows 10. Les anciennes applications ainsi que leurs configurations sont transférées sur le nouveau PC si possible, sinon une mise à jour est faite vers un applicatif plus récent. Le PC est ensuite testé et validé dans un environnement de test puis mis en production.
        \newline
        \includegraphics[width=420px]{../Images/migration.png}
        \item Formation des utilisateurs aux nouvelles technologies mises en place au sein de l'entreprise.
    \end{itemize}
    \newpage

    \section{Cahier des charges}
        L’objectif de ce projet est de pouvoir mesurer la qualité de l’air des usines ainsi que d’envoyer des
    alertes si des seuils sont dépassés.
    Pour cela, un capteur de qualité d’air sera relié à un microcontrôleur qui va se charger de récupérer
    les données. Il va ensuite les chiffrer via SSL et les envoyer au broker MQTT hébergé sur le
    serveur Linux.
    Le serveur Linux va ensuite récupérer les informations, les stocker dans une base de données, les
    afficher sur une Dashboard et envoyer des alertes si besoin.

    
    \subsection{Diagramme de contexte :}
    Le capteur de QAI effectue des mesures à intervalle régulier, ces mesures sont transmises au
    microcontrôleur STM32 qui les transmet à son tour au serveur. Le serveur traite les données, les
    stocke et les affiche sur une Dashboard. La Dashboard est accessible par les utilisateurs pour
    visualiser les données et par les administrateurs pour configurer le système :
    \includegraphics[width=420px]{../Images/Contexte_CDC.png}
    \newpage

    \subsection{Diagramme UML de cas d’utilisation :}
    Le technicien informatique accède à l’interface de paramétrage du système après c’être authentifié.
    Il peut paramétrer les alertes et le système de capteurs depuis cette interface.
    L’utilisateur peut consulter les données après authentification et reçoit des alertes mail si un seuil
    de QAI est dépassé.
    Le microcontrôleur récupère les données des capteurs et les transmet au serveur de centralisation.
    \newline
    \includegraphics[width=420px]{../Images/Cas_utilisation.png}
    \newpage

    \subsection{Diagramme de séquence de visualisation des données :}
    Lorsqu’un client souhaite visualiser les données des capteurs, il va accéder à la Dashboard via un
    navigateur internet. L’URL sera convertie en adresse IP par le serveur DNS, puis le serveur
    Authelia se chargera de l’authentification de l’utilisateur. Enfin, Node-Red va récupérer les données
    dans la base de données MariaDB et les afficher sur le poste client
    \includegraphics[width=420px]{../Images/Sequence_visu.png}
    \newpage

    \subsection{Diagramme de séquence de réception / envoi des données :}
    Le capteur de QAI transmet les données récoltées au microcontrôleur STM32 qui va ensuite les
    chiffrer à l’aide de SSL/TLS puis les envoyer au broker MQTT. Le broker MQTT déchiffre les
    données et les transmet aux clients qui ont souscrits au topic. Le script Python est un de ces clients
    et il se charge de récupérer les données et de les écrire dans la base de données
    \newline
    \includegraphics[width=420px]{../Images/Sequence_envoie.png}

    \subsection{Diagramme de Merise :}
    Chaque capteur est identifié par un numéro unique ainsi qu’un nom donné par l’administrateur.
    Le capteur mesure plusieurs types de grandeurs physiques, spécifiées dans id type.
    Chaque mesure est liée au capteur, au type de mesure et horodatée
    \newline
    \includegraphics[width=420px]{../Images/merise.png}
    \newpage
    
    \subsection{Diagramme de déploiement :}
    Le système sera déployé en 4 parties distinctes :
    Un ou plusieurs sous-ensembles capteurs + microcontrôleurs seront installés dans les usines.
    Ils communiqueront avec le serveur GNU/Linux qui a un broker MQTT, un serveur Node-Red, une
    base de données MariaDB, un serveur mail ainsi que divers services de sécurisation.
    Ce serveur sera également utilisé par les PC clients des collaborateurs pour visualiser les données
    ainsi que par les PC clients de l’IT pour visualiser les données et configurer le système. Tous ces
    sous-ensembles communiqueront par le réseau.
    \newline
    \includegraphics[width=420px]{../Images/Deploiement.png}

    \newpage
    \subsection{Contraintes techniques et économiques}


    \subsubsection{Contraintes Techniques}
    Le microcontrôleur devra être programmé en Rust ou en C.\newline
    La distribution Linux devra être Debian ou Arch Linux.\newline
    Les conteneurs doivent respecter la norme OCI.\newline
    Les conteneurs doivent être gérés avec Kubernetes.\newline
    La Dashboard sera faite avec Node-Red.\newline
    Un serveur DNS doit être fonctionnel pour permettre la résolution de l’url de la dashboard.\newline
    La connexion au serveur doit se faire via un protocole sécurisé.\newline
    Le serveur doit être sécurisé avec un pare-feu.\newline
    Les données doivent transiter via le protocole MQTT sécurisé avec SSL/TLS.\newline
    Le capteur doit être utilisable en environnement industriel (résistant à la poussière et à la chaleur).\newline
    Le système doit être sauvegardé pour permettre une restauration complète du système en cas de panne ou de cyberattaque.\newline
    Le code doit être versionné avec l’outil git et documenté par le biais d’un wiki.\newline

    \subsubsection{Contraintes Economiques}
    Achat Microcontrôleur DevBoard : 30€ \newline
    Achat Capteur : 20€

    \subsection{Ressources Utilisées}
    \begin{tabular}{| C{30px} | C{5cm} | C{8cm} |}
        \hline
        \includegraphics[width=30px]{../Images/buildah.png} & Buildah & Création des conteneurs OCI \\
        \hline
        \includegraphics[width=30px]{../Images/podman.png} & Podman  & Test des conteneurs OCI \\
        \hline
        \includegraphics[width=30px]{../Images/kubernetes.png} & Kubernetes  & Gestion des conteneurs OCI \\
        \hline
        \includegraphics[width=30px]{../Images/node-red.png} & Node-Red & Création de la dashboard \\
        \hline
        \includegraphics[width=30px]{../Images/authelia.png} & Authelia & Gestion des utilisateurs Node-Red \\
        \hline
        \includegraphics[width=30px]{../Images/bind.png} & Bind 9 & Serveur DNS \\
        \hline
        \includegraphics[width=30px]{../Images/ssh.png} & SSH & Administration à distance sécurisée \\
        \hline
        \includegraphics[width=30px]{../Images/postfix.png} & PostFix & Serveur SMTP pour l'envoi des mails d'alerte \\
        \hline
        \includegraphics[width=30px]{../Images/ufw.png} & UFW & Sécurisation réseau du serveur (pare-feux) \\
        \hline
        \includegraphics[width=30px]{../Images/mqtt_logo.png} & MQTT & Récupération / Envoie des données du capteur \\
        \hline
        \includegraphics[width=30px]{../Images/espressif.png} & STM32H753ZI / ESP32-Wroom32E & Traitement des données du capteur et envoie au serveur \\
        \hline
        \includegraphics[width=30px]{../Images/rust.png} & C / Rust & Programmation du STM32 / ESP32 \\
        \hline
        \includegraphics[width=30px]{../Images/nuc.jpg} & PC Intel NUC & Serveur Linux \\
        \hline
        \includegraphics[width=30px]{../Images/sds011.jpg} & PMS5003 / SDS011 / MQ135 / ADA4632 & Capteur de QAI \\
        \hline
        \includegraphics[width=30px]{../Images/mariadb.png} & MariaDB & Base de données \\
        \hline
        \includegraphics[width=30px]{../Images/debian.png} & Debian / Arch / Alpine Linux / Ubuntu & Système d'exploitation \\
        \hline
        \includegraphics[width=30px]{../Images/veeam.png} & Timeshift / Veeam & Backup du système et des conteneurs OCI \\
        \hline
        \includegraphics[width=30px]{../Images/python.png} & Python3 / Rust & Script de récupération et d'envoie des données depuis MQTT vers MariaDB \\
        \hline
        \includegraphics[width=30px]{../Images/sphinx.png} & Sphinx / RustDoc & Création du wiki documentant le code \\
        \hline
        \includegraphics[width=30px]{../Images/git.png} & Git + Github & Versionning du code \\
        \hline
    \end{tabular}
    \newpage
    \section{Firmware}
    \section{Configuration}
    \section{Conclusion}
    \section{Annexe}
    Documentations liées au système
    \subsection{Mosquitto}
    \includegraphics[width=60px]{../Images/mqtt_logo.png}
    \newline
    \href{https://mosquitto.org/documentation/}{Documentation MQTT}
    \newline
    Schéma de fonctionnement du protocol MQTT :
    \newline
    \includegraphics[width=420px]{../Images/MQTT.png}
    \subsection{MariaDB}
    \includegraphics[width=60px]{../Images/mariadb.png}
    \newline
    \href{https://mariadb.org/documentation/}{Documentation MariaDB}
    \subsection{ESP32}
    \subsection{Debian}
    \includegraphics[width=60px]{../Images/debian.png}
    \newline
    \href{https://www.debian.org/doc/}{Documentation Debian}

\end{document}
